% Modelo Beamer UFC original de Maurício Moreira Neto.
\ifdefined\Resumo\PassOptionsToClass{handout}{beamer}\fi
\documentclass{libs/ufc_format}
\usepackage{amsmath,amssymb,booktabs,tikz,circuitikz}
\usetikzlibrary{arrows.meta,positioning,calc,shapes.geometric}
% A geometria 4:3 e todos os elementos do tema vêm da classe ORIGINAL.
% Não substituir fundo, logo, cabeçalho, rodapé, fontes ou blocos do template.
\hypersetup{pdflang=pt-BR}
\usepackage{adjustbox}
\newcommand{\VDD}{V_{\mathrm{DD}}}
\newcommand{\GND}{\mathrm{GND}}
\newcommand{\fonte}[1]{\par\vspace{1mm}{\tiny\color{blue_theme}#1\par}}
\newcommand{\ideia}[1]{\begin{block}{Ideia-chave}#1\end{block}}
\newcommand{\pergunta}[1]{\begin{exampleblock}{Para discutir}#1\end{exampleblock}}
\newcommand{\duas}[2]{\begin{columns}[T,onlytextwidth]\begin{column}{.48\textwidth}#1\end{column}\begin{column}{.48\textwidth}#2\end{column}\end{columns}}
\newcommand{\topico}[2]{\begin{block}{#1}#2\end{block}}
\newcommand{\capa}[2]{\subtitle{#1\\[1mm]{\small #2}}\begin{frame}{}\maketitle\end{frame}}
% Redimensiona somente diagramas e tabelas que excedam a largura de texto.
\tikzset{caixa/.style={draw=blue_theme,fill=blue_theme!5,rounded corners=2pt,align=center,minimum height=9mm,font=\small},seta/.style={-{Latex},thick,blue_theme}}
\ctikzset{transistors/scale=.78}
\newcommand{\inversor}{%
\begin{circuitikz}[scale=.85,transform shape]
\node[pmos] (p) at (0,1.3) {};\node[nmos] (n) at (0,-1.1) {};
\draw (p.S) -- (0,2.2) node[above]{$\VDD$};
\draw (p.D) -- (n.D);\draw (0,.1) to[short,*-] (1.35,.1) node[right]{$Y$};
\draw (n.S) -- (0,-2) node[ground]{};
\draw (p.G) -- (-1.4,1.3) -- (-1.4,-1.1) -- (n.G);
\draw (-1.4,.1) to[short,*-] (-2.05,.1) node[left]{$A$};
\node[right] at (.25,1.3) {PMOS};\node[right] at (.25,-1.1) {NMOS};
\end{circuitikz}}
\newcommand{\nandcircuito}{%
\begin{circuitikz}[scale=.68,transform shape]
\node[pmos] (pa) at (-1.2,1.7) {};\node[pmos] (pb) at (1.4,1.7) {};
\node[nmos] (na) at (0,-.7) {};\node[nmos] (nb) at (0,-2.3) {};
\draw (pa.S) -- (-1.2,2.6) -- (1.4,2.6) -- (pb.S);\node[above] at (.1,2.6) {$\VDD$};
\draw (pa.D) -- (-1.2,.5) -- (1.4,.5) -- (pb.D);
\draw (0,.5) -- (na.D);\draw (na.S) -- (nb.D);\draw (nb.S) -- (0,-3.15) node[ground]{};
\draw (1.4,.5) to[short,*-] (2.15,.5) node[right]{$Y$};
\draw (pa.G) node[left]{$A$};\draw (pb.G) node[left]{$B$};\draw (na.G) node[left]{$A$};\draw (nb.G) node[left]{$B$};
\end{circuitikz}}
\newcommand{\norcircuito}{%
\begin{circuitikz}[scale=.68,transform shape]
\node[pmos] (pa) at (0,2) {};\node[pmos] (pb) at (0,.4) {};
\node[nmos] (na) at (-1.2,-1.8) {};\node[nmos] (nb) at (1.4,-1.8) {};
\draw (pa.S) -- (0,2.85) node[above]{$\VDD$};\draw (pa.D) -- (pb.S);\draw (pb.D) -- (0,-.6);
\draw (na.D) -- (-1.2,-.6) -- (1.4,-.6) -- (nb.D);
\draw (1.4,-.6) to[short,*-] (2.15,-.6) node[right]{$Y$};
\draw (na.S) -- (-1.2,-2.8) -- (1.4,-2.8) -- (nb.S);\draw (0,-2.8) node[ground]{};
\draw (pa.G) node[left]{$A$};\draw (pb.G) node[left]{$B$};\draw (na.G) node[left]{$A$};\draw (nb.G) node[left]{$B$};
\end{circuitikz}}
\newcommand{\aoi}{%
\begin{circuitikz}[scale=.64,transform shape]
\node[pmos] (pa) at (-1.3,2) {};\node[pmos] (pb) at (1.3,2) {};\node[pmos] (pc) at (0,.4) {};
\draw (pa.S) -- (-1.3,2.8) -- (1.3,2.8) -- (pb.S);\node[above] at (0,2.8) {$\VDD$};
\draw (pa.D) -- (-1.3,1.2) -- (1.3,1.2) -- (pb.D);\draw (0,1.2) -- (pc.S);
\draw (pc.D) -- (0,-.4) -- (2.2,-.4) node[right]{$Y$};
\node[nmos] (na) at (-1.3,-1.3) {};\node[nmos] (nb) at (-1.3,-2.9) {};\node[nmos] (nc) at (1.3,-2.1) {};
\draw (na.D) -- (-1.3,-.4) -- (1.3,-.4) -- (nc.D);\draw (na.S) -- (nb.D);
\draw (nb.S) -- (-1.3,-3.75) -- (1.3,-3.75) -- (nc.S);\draw (0,-3.75) node[ground]{};
\draw (pa.G) node[left]{$A$};\draw (pb.G) node[left]{$B$};\draw (pc.G) node[left]{$C$};
\draw (na.G) node[left]{$A$};\draw (nb.G) node[left]{$B$};\draw (nc.G) node[left]{$C$};
\end{circuitikz}}
% Figuras vetoriais autorais: parâmetros controlam etapas da apresentação.
\definecolor{eletron}{RGB}{0,83,134}
\definecolor{lacuna}{RGB}{199,123,0}
\definecolor{ativo}{RGB}{0,132,75}
\tikzset{fio/.style={line width=.9pt,black!65},ligado/.style={line width=2pt,ativo},cortado/.style={line width=.7pt,black!25},anotacao/.style={font=\small,align=center},passo/.style={font=\scriptsize,text=blue_theme}}
\newcommand{\chamada}[1]{\par\medskip\centering\textcolor{blue_theme}{\textbf{#1}}\par}
\newcommand{\etapa}[2]{\par\smallskip{\scriptsize\color{blue_theme}Etapa #1\enspace•\enspace #2}\par}
% PN: 0 contato inicial; 1 difusão; 2 equilíbrio; 3 direta; 4 reversa.
\newcommand{\juncaoPN}[1]{%
\begin{tikzpicture}[x=.87cm,y=.85cm,font=\small]
\path[use as bounding box] (-.3,-1.65) rectangle (10.3,4.9);
\fill[lacuna!9] (0,0) rectangle (5,3);\fill[eletron!9] (5,0) rectangle (10,3);
\node[blue_theme] at (2,3.5){\textbf{Região P}};\node[blue_theme] at (8,3.5){\textbf{Região N}};
\ifnum#1>1
 \def\depl{1.25}\ifnum#1=3\def\depl{.6}\fi\ifnum#1=4\def\depl{2}\fi
 \fill[gray!25] (5-\depl,0) rectangle (5+\depl,3);
 \draw[densely dashed,gray] (5-\depl,0)--(5-\depl,3);\draw[densely dashed,gray] (5+\depl,0)--(5+\depl,3);
 \foreach \y in {.55,1.25,1.95,2.6}{\node[eletron] at (5-.5*\depl,\y){$-$};\node[lacuna] at (5+.5*\depl,\y){$+$};}
 \draw[<->,blue_theme] (5-\depl,4.15)--node[above]{\strut Depleção}(5+\depl,4.15);
 \draw[-{Latex},very thick,blue_theme] (5+\depl,-.55)--node[below]{Campo interno}(5-\depl,-.55);
\fi
\foreach \x/\y in {.7/.6,1.6/1.4,2.5/.65,.8/2.3,2.6/2.4}{\draw[lacuna,fill=white,line width=1pt] (\x,\y) circle(.13);}
\foreach \x/\y in {7.5/.6,8.4/1.4,9.3/.65,7.4/2.4,9.2/2.3}{\fill[eletron] (\x,\y) circle(.12);}
\ifnum#1<2
 \foreach \x/\y in {3.6/.8,4.3/2,3.4/2.5}{\draw[lacuna,fill=white,line width=1pt] (\x,\y) circle(.13);}
 \foreach \x/\y in {5.7/.8,6.4/2,6.6/2.5}{\fill[eletron] (\x,\y) circle(.12);}
\fi
\ifnum#1=1
 \draw[-{Latex},very thick,eletron] (6.3,1.2)--node[above]{elétrons}(3.8,1.2);
 \draw[-{Latex},very thick,lacuna] (3.6,2.3)--node[above]{lacunas}(6.1,2.3);
 \node[font=\scriptsize] at (5,-1.25){Difusão e recombinação próximas à interface};
\fi
\ifnum#1=0\node at (5,-1){\textcolor{lacuna}{$\circ$} lacuna \qquad\textcolor{eletron}{$\bullet$} elétron};\fi
\ifnum#1=2\node[font=\scriptsize] at (5,-1.3){Sinais na faixa cinza: íons fixos, não portadores móveis};\fi
\ifnum#1=3
 \node at (1.2,4.3){$+$};\node at (8.8,4.3){$-$};
 \node[font=\scriptsize] at (5,-1.3){Direta: barreira menor, corrente maior};
\fi
\ifnum#1=4
 \node at (1.2,4.3){$-$};\node at (8.8,4.3){$+$};
 \node[font=\scriptsize] at (5,-1.3){Reversa: barreira maior, pequena fuga antes da ruptura};
\fi
\end{tikzpicture}}
% NMOS planar: 0 sem canal forte; 1 depleção; 2 canal; 3 corrente.
\newcommand{\moscanal}[1]{%
\begin{tikzpicture}[x=.85cm,y=.8cm,font=\small]
\path[use as bounding box] (-.3,-1.1) rectangle (11,5.2);
\fill[eletron!8,draw=blue_theme] (0,0) rectangle (10,2.8);
\fill[yellow!35,draw=blue_theme] (.5,1.7) rectangle (2.5,2.8);
\fill[yellow!35,draw=blue_theme] (7.5,1.7) rectangle (9.5,2.8);
\node at (1.5,2.2){$n^+$};\node at (8.5,2.2){$n^+$};
\fill[gray!40,draw=blue_theme] (2.7,2.85) rectangle (7.3,3.08);
\fill[eletron!55,draw=blue_theme] (2.7,3.08) rectangle (7.3,3.55);
\node at (5,.5){Corpo P};\node at (1.5,3.9){Fonte};\node at (8.5,3.9){Dreno};\node at (5,4){Porta};
\draw (1.5,2.8)--(1.5,3.6);\draw (8.5,2.8)--(8.5,3.6);\draw (5,3.55)--(5,3.8);
\node[font=\scriptsize,anchor=west] at (8.1,4.72){Dielétrico};\draw[-{Latex},black!50] (8.1,4.35)--(7.65,4.35)--(7.25,2.97);
\ifnum#1<1
 \foreach \x/\y in {3.3/1.2,4.3/2.1,5.2/1.4,6.5/2.1}{\draw[lacuna,fill=white] (\x,\y) circle(.12);}
 \node at (5,4.75){$V_{GS}$ baixo};
\else
 \fill[gray!20] (2.5,1.55) rectangle (7.5,2.8);
 \foreach \x in {3,4,5,6,7}{\node[eletron,font=\scriptsize] at (\x,2.1){$-$};}
 \foreach \x in {3.2,4.4,5.6,6.8}{\node[red] at (\x,3.32){$+$};}
 \ifnum#1=1
  \node at (5,4.75){A tensão de porta aumenta};
  \draw[-{Latex},lacuna,thick] (4,1.55)--(4,.8);\draw[-{Latex},lacuna,thick] (6,1.55)--(6,.8);
  \node[font=\scriptsize] at (5,-.65){Lacunas se afastam da superfície: depleção};
 \else
  \fill[eletron!70] (2.5,2.57) rectangle (7.5,2.8);
  \foreach \x in {2.9,3.6,4.3,5,5.7,6.4,7.1}{\node[white,font=\scriptsize] at (\x,2.69){$-$};}
  \node at (5,4.75){$V_{GS}>V_{TN}$: canal de elétrons};
  \ifnum#1=3
   \draw[-{Latex},white,line width=3.5pt] (8.4,2.69)--(1.6,2.69);\draw[-{Latex},line width=2pt,ativo] (8.4,2.69)--(1.6,2.69);\node[ativo] at (5,1.15){Corrente convencional $I_D$};
   \node[font=\scriptsize] at (5,-.65){Com $V_D>V_S$, há corrente pelo canal};
  \else\node[font=\scriptsize] at (5,-.65){Inversão: a superfície conduz por elétrons};\fi
 \fi
\fi
\end{tikzpicture}}
% Circuitos por estado: caminho ativo verde; transistor bloqueado cinza.
\newcommand{\inversorestado}[1]{%
\begin{circuitikz}[scale=.9,transform shape]
\path[use as bounding box] (-3,-2.9) rectangle (3,3);
\ifnum#1=0\def\pstyle{ativo}\def\nstyle{black!25}\def\saida{1}\else\def\pstyle{black!25}\def\nstyle{ativo}\def\saida{0}\fi
\node[pmos,color=\pstyle] (p) at (0,1.3) {};\node[nmos,color=\nstyle] (n) at (0,-1.1) {};
\draw[ligado,color=\pstyle] (p.S)--(0,2.2) node[above,black]{$\VDD$};
\draw[ligado,color=\pstyle] (p.D)--(0,.1);
\draw[ligado,color=\nstyle] (n.D)--(0,.1);\draw[ligado,color=\nstyle] (n.S)--(0,-2.1) node[ground,black]{};
\draw[fio] (p.G)--(-1.5,1.3)--(-1.5,-1.1)--(n.G);\draw[fio] (-1.5,.1)--(-2.1,.1) node[left]{$A=#1$};
\draw[fio] (0,.1) to[short,*-] (1.3,.1) node[right]{$Y=\saida$};
\node[font=\small,text=\pstyle,right] at (.5,1.3){PMOS};\node[font=\small,text=\nstyle,right] at (.5,-1.1){NMOS};
\node[font=\scriptsize,ativo] at (0,-2.8){Verde: caminho que fixa a saída};
\end{circuitikz}}
\newcommand{\portaestado}[3]{%
\begin{circuitikz}[scale=.78,transform shape]
\path[use as bounding box] (-3,-3.8) rectangle (3.7,3.6);
\ifnum#2=0\def\pa{ativo}\def\na{black!25}\else\def\pa{black!25}\def\na{ativo}\fi
\ifnum#3=0\def\pb{ativo}\def\nb{black!25}\else\def\pb{black!25}\def\nb{ativo}\fi
\ifnum#1=0
 \pgfmathtruncatemacro{\out}{1-#2*#3}
 \node[pmos,color=\pa] (pa) at (-1.2,1.7) {};\node[pmos,color=\pb] (pb) at (1.4,1.7) {};
 \node[nmos,color=\na] (na) at (0,-.7) {};\node[nmos,color=\nb] (nb) at (0,-2.3) {};
 \draw[fio] (pa.S)--(-1.2,2.6)--(1.4,2.6)--(pb.S);\node[above] at (.1,2.6){$\VDD$};
 \draw[ligado,color=\pa] (pa.D)--(-1.2,.5)--(0,.5);\draw[ligado,color=\pb] (pb.D)--(1.4,.5)--(0,.5);
 \draw[ligado,color=\na] (0,.5)--(na.D);\draw[ligado,color=\na] (na.S)--(0,-1.5);
 \draw[ligado,color=\nb] (0,-1.5)--(nb.D);\draw[ligado,color=\nb] (nb.S)--(0,-3.2) node[ground,black]{};
 \draw[fio] (0,.5) to[short,*-] (2.2,.5) node[right]{$Y=\out$};
\else
 \pgfmathtruncatemacro{\out}{(1-#2)*(1-#3)}
 \node[pmos,color=\pa] (pa) at (0,2) {};\node[pmos,color=\pb] (pb) at (0,.4) {};
 \node[nmos,color=\na] (na) at (-1.2,-1.8) {};\node[nmos,color=\nb] (nb) at (1.4,-1.8) {};
 \draw[ligado,color=\pa] (pa.S)--(0,2.85) node[above,black]{$\VDD$};\draw[ligado,color=\pa] (pa.D)--(0,1.2);
 \draw[ligado,color=\pb] (0,1.2)--(pb.S);\draw[ligado,color=\pb] (pb.D)--(0,-.6);
 \draw[ligado,color=\na] (na.D)--(-1.2,-.6)--(0,-.6);\draw[ligado,color=\nb] (nb.D)--(1.4,-.6)--(0,-.6);
 \draw[ligado,color=\na] (na.S)--(-1.2,-2.8);\draw[ligado,color=\nb] (nb.S)--(1.4,-2.8);
 \draw[fio] (-1.2,-2.8)--(1.4,-2.8);\draw (0,-2.8) node[ground]{};
 \draw[fio] (0,-.6) to[short,*-] (2.2,-.6) node[right]{$Y=\out$};
\fi
\draw (pa.G) node[left,black]{$A=#2$};\draw (pb.G) node[left,black]{$B=#3$};
\draw (na.G) node[left,black]{$A=#2$};\draw (nb.G) node[left,black]{$B=#3$};
\end{circuitikz}}

% Figuras originais para a abertura histórica. Referências: dados/FONTES_CONTEXTO.md.
% As cores distinguem Gate, canal e interconexões; desenhos sem escala.
\newcommand{\linhaHistorica}{%
\begin{tikzpicture}[x=1cm,y=1cm,font=\small]
\draw[seta] (0,0)--(10.3,0);
\foreach \x/\ano/\titulo/\sub in {0.4/1947/Transistor/Bell Labs,2.8/1958--59/Circuito integrado/Kilby e Noyce,5.2/1960--63/MOS e CMOS/Controle por campo,7.6/1971/Microprocessador/Intel 4004,10/Hoje/Sistemas integrados/Múltiplas funções}{%
\fill[blue_theme] (\x,0) circle(.09);
\node[blue_theme,above=3mm,align=center,font=\small\bfseries] at (\x,0){\ano};
\node[below=3mm,align=center,text width=2.25cm,font=\scriptsize] at (\x,0){\textbf{\titulo}\\[1mm]\sub};}
% Dispositivo discreto, circuito, CPU, SoC: motivos conceituais.
\draw[blue_theme,thick] (.05,1.5) rectangle (.75,2.25);
\foreach \x in {.15,.4,.65}{\draw[blue_theme] (\x,1.5)--(\x,1.1);}
\draw[blue_theme,thick] (2,1.2) rectangle (3.6,2.5);
\foreach \x in {2.2,2.7,3.2}{\draw[blue_theme] (\x,1.6)--++(.2,.45)--++(.2,-.45);}
\draw[blue_theme,thick] (4.6,1.4)--(5.8,1.4);\draw[blue_theme,thick] (4.6,1.65)--(5.8,1.65);\draw[blue_theme,line width=4pt] (4.9,2)--(5.5,2);
\draw[blue_theme,thick] (6.8,1.2) rectangle (8.4,2.5);\node at (7.6,1.85){CPU};
\draw[blue_theme,thick] (9.1,1.2) rectangle (10.9,2.5);
\foreach \x/\y in {9.25/1.35,10.05/1.35,9.25/1.95,10.05/1.95}{\fill[blue_theme!20] (\x,\y) rectangle ++(.65,.4);}
\path[use as bounding box] (-.8,-1.3) rectangle (11.2,2.9);
\end{tikzpicture}}

\newcommand{\graficoMoore}{%
\begin{tikzpicture}[x=.86cm,y=.19cm,font=\small]
\draw[->,blue_theme] (0,0)--(9,0) node[right]{anos};
\draw[->,blue_theme] (0,0)--(0,19) node[above,align=center]{quantidade\\relativa};
\foreach \y in {1,2,4,8,16}{\draw[gray!25] (0,\y)--(8,\y);\node[left] at (0,\y){\y};}
\foreach \x in {0,2,4,6,8}{\draw (\x,0)--(\x,-.4) node[below]{\x};}
\draw[blue_theme,very thick,domain=0:8,samples=60] plot (\x,{2^(\x/2)});
\foreach \x/\y in {0/1,2/2,4/4,6/8,8/16}{\fill[blue_theme] (\x,\y) circle[x radius=.07,y radius=.32];}
\node[align=center,fill=white,inner sep=2mm] at (3.6,14){\textbf{Dobrar a cada dois anos}\\$N(t)=N_0\,2^{t/2}$};
\path[use as bounding box] (-1.4,-3) rectangle (10.2,24);
\end{tikzpicture}}

\newcommand{\graficoProcessos}{%
\begin{tikzpicture}[x=.175cm,y=.79cm,font=\scriptsize]
\fill[blue_theme!6] (47,-.1) rectangle (56,4.8);
\draw[->,blue_theme] (0,0)--(57,0);\draw[->,blue_theme] (0,0)--(0,4.8);
\node[anchor=west] at (0,5.15){Referência de processo (nm) · eixo logarítmico};
\foreach \y/\t in {0/1,1/10,2/100,3/{1\,000},4/{10\,000}}{\draw[gray!25] (0,\y)--(56,\y);\node[left] at (0,\y){\t};}
\draw[blue_theme,very thick] (0,4)--(7,3.4771)--(14,3.1761)--(22,2.9031)--(29,2.2553)--(37,1.6532);
\foreach \x/\y/\valor/\ano in {0/4/{10\,000}/1971,14/3.1761/{1\,500}/1985,29/2.2553/180/2000,37/1.6532/45/2008}{%
\fill[blue_theme] (\x,\y) circle[x radius=.45,y radius=.1];
\ifnum\x>0\node[above=1.5mm] at (\x,\y){\valor};\fi\node[below=2mm] at (\x,0){\ano};}
\foreach \x/\y/\nome/\ano in {49/.699/N5/2020,54/.301/N2/2025}{%
\draw[lacuna,thick,fill=white] (\x,\y) circle[x radius=.5,y radius=.11];
\node[above=2mm,lacuna] at (\x,\y){\nome};\node[below=2mm] at (\x,0){\ano};}
\node[align=center,text width=1.7cm] at (51.5,3.4){Nós comerciais\\recentes};
\node[anchor=west] at (0,-1.1){$1\,\mu$m = $1\,000$ nm};
\path[use as bounding box] (-7,-1.5) rectangle (58,5.6);
\end{tikzpicture}}

\newcommand{\litografiaEtapa}[1]{%
\begin{tikzpicture}[x=1cm,y=1cm,font=\small]
\fill[blue_theme!15] (0,0) rectangle (7,.7);\node at (3.5,.32){substrato};
\ifnum#1<3
\fill[blue_theme!65] (0,.7) rectangle (7,.95);
\else
\foreach \x in {0,2.8,5.6}{\fill[blue_theme!65] (\x,.7) rectangle ++(1.4,.25);}
\fi
\ifnum#1=1
\fill[yellow!70] (0,.95) rectangle (7,1.3);
\foreach \x in {1.4,4.2}{\fill[red!12] (\x,.95) rectangle ++(1.4,.35);\fill[red!10] (\x,1.4) rectangle ++(1.4,2.2);\draw[-{Latex},red] (\x+.7,3.5)--(\x+.7,1.45);}
\draw[blue_theme,line width=4pt] (0,3.8)--(1.4,3.8) (2.8,3.8)--(4.2,3.8) (5.6,3.8)--(7,3.8);
\node[left] at (0,3.8){máscara};\node[right,align=left,font=\scriptsize] at (7.2,2.6){Luz projetada\\sobre o resiste};
\else
\foreach \x in {0,2.8,5.6}{\fill[yellow!70] (\x,.95) rectangle ++(1.4,.35);}
\ifnum#1=2
\foreach \x in {2.1,4.9}{\draw[seta] (\x,2.6)--(\x,1.5);}
\node[align=center] at (3.5,3.5){Revelação remove as partes expostas\\de um resiste positivo.};
\else
\foreach \x in {2.1,4.9}{\draw[-{Latex},thick,red] (\x,2.6)--(\x,.85);}
\node[align=center] at (3.5,3.5){Gravação transfere o padrão\\à camada abaixo do resiste.};
\fi\fi
\draw[blue_theme] (0,0) rectangle (7,.7);
\draw[-,blue_theme] (7,.82)--(7.6,.82) node[right,font=\scriptsize]{camada a padronizar};
\draw[-,lacuna] (6.4,1.2)--(7.6,1.7) node[right,font=\scriptsize]{resiste};
\path[use as bounding box] (-1.5,-.4) rectangle (10.1,4.5);
\end{tikzpicture}}

\newcommand{\transistoresModernos}{%
\begin{tikzpicture}[x=1cm,y=1cm,font=\small]
\foreach \x/\nome in {0/Planar,3.8/FinFET,7.6/{GAA · nanosheets}}{%
\node[blue_theme,font=\small\bfseries] at (\x+1.25,3.2){\nome};
\fill[gray!12] (\x,0) rectangle ++(2.5,.5);}
% Cortes transversais ao canal. Dielétrico branco entre Gate e silício.
\fill[ativo!60] (.4,.5) rectangle (2.1,.85);
\fill[blue_theme!70] (.35,1) rectangle (2.15,1.5);
\fill[blue_theme!70] (4.35,.55) rectangle (5.75,2.4);
\fill[white] (4.75,.5) rectangle (5.35,2.05);
\fill[ativo!60] (4.87,.5) rectangle (5.23,1.93);
\fill[blue_theme!70] (8.1,.65) rectangle (9.6,2.55);
\foreach \y in {.9,1.5,2.1}{\fill[white,rounded corners=2pt] (8.3,\y) rectangle ++(1.1,.3);\fill[ativo!60,rounded corners=1pt] (8.4,\y+.08) rectangle ++(.9,.14);}
\node[align=center,font=\scriptsize] at (1.25,-.5){Gate sobre\\a superfície};
\node[align=center,font=\scriptsize] at (5.05,-.5){Gate envolve\\três faces};
\node[align=center,font=\scriptsize] at (8.85,-.5){Gate envolve\\cada canal};
\fill[blue_theme!70] (2,-1.45) rectangle ++(.25,.2);\node[right,font=\scriptsize] at (2.3,-1.35){Gate};
\fill[ativo!60] (5,-1.45) rectangle ++(.25,.2);\node[right,font=\scriptsize] at (5.3,-1.35){canal};
\path[use as bounding box] (-.3,-1.7) rectangle (10.4,3.6);
\end{tikzpicture}}

\newcommand{\pilhaCMOSdois}{%
\begin{tikzpicture}[x=1cm,y=1cm,font=\small]
\foreach \y/\cor/\titulo in {0/gray/Alimentação e acesso externo,1.2/ativo/Lógica densa,2.4/blue_theme/Memória,3.6/lacuna/Interfaces e funções especializadas}{%
\fill[\cor!14] (0,\y)--(5,\y)--(5.8,\y+.45)--(.8,\y+.45)--cycle;
\draw[\cor,thick] (0,\y)--(5,\y)--(5.8,\y+.45)--(.8,\y+.45)--cycle;
\node[anchor=west,font=\scriptsize] at (6,\y+.2){\titulo};}
\foreach \x in {1,2,3,4}{\draw[blue_theme,thick,densely dashed] (\x,.5)--(\x,3.5);}
\node[anchor=west,font=\scriptsize,align=left] at (6,-.9){Conexões 3D densas\\entre camadas funcionais};
\path[use as bounding box] (-.3,-1.4) rectangle (10.6,4.5);
\end{tikzpicture}}

% Diagramas didáticos originais. Sem parâmetros de um PDK.
\newcommand{\corteCMOS}[1]{%
\begin{tikzpicture}[x=.88cm,y=.88cm,font=\small]
\path[use as bounding box] (-.4,-.7) rectangle (11.4,5.4);
\fill[blue_theme!10] (0,0) rectangle (11,1.4);
\fill[lacuna!17] (.55,.2) rectangle (5.05,1.4);
\draw[lacuna,thick] (.55,1.4)--(.55,.2)--(5.05,.2)--(5.05,1.4);
\fill[gray!35] (5.35,.9) rectangle (6.45,1.4);
\node[font=\scriptsize] at (2.8,.5){poço N};
\node[font=\scriptsize] at (8.1,.45){substrato P};
\node[blue_theme,font=\bfseries] at (2.8,-.35){PMOS};
\node[blue_theme,font=\bfseries] at (8.55,-.35){NMOS};
\ifnum#1>1
\fill[white] (0,1.4) rectangle (11,1.54);
\foreach \x/\tipo/\cor in {1.1/{p$^+$}/lacuna,3.5/{p$^+$}/lacuna,7/{n$^+$}/eletron,9.4/{n$^+$}/eletron}{%
\fill[\cor!40] (\x,1) rectangle ++(.9,.4);\node[font=\scriptsize] at (\x+.45,1.18){\tipo};}
\foreach \x in {2.2,8.15}{\fill[lacuna!65] (\x,1.43) rectangle ++(.9,.1);\fill[blue_theme!65] (\x,1.59) rectangle ++(.9,.38);}
\ifnum#1=2
\draw[blue_theme] (2.65,2.05)--(2.65,2.6) node[above,font=\scriptsize]{porta};
\draw[lacuna] (8.6,1.48)--(9.5,2.5) node[above,font=\scriptsize]{dielétrico};
\node[font=\scriptsize,align=center] at (5.5,3.9){Os dois dispositivos têm regiões de fonte e dreno\\e uma porta isolada do semicondutor.};
\fi
\else
\node[font=\scriptsize,align=center] at (5.5,3.5){Regiões com dopagens diferentes\\permitem construir os dispositivos complementares.};
\draw[blue_theme] (5.9,1.4)--(5.9,2.15) node[above,font=\scriptsize]{isolação};
\fi
\ifnum#1=3
\foreach \x in {1.55,3.95,7.45,9.85}{\fill[black!55] (\x-.075,1.4) rectangle (\x+.075,2.55);}
\draw[line width=2pt,ativo] (1.55,2.55)--(1.55,3.55) node[above]{$\VDD$};
\draw[line width=2pt,ativo] (9.85,2.55)--(9.85,3.55) node[above]{GND};
\draw[line width=2pt,blue_theme] (3.95,2.55)--(3.95,3)--(7.45,3)--(7.45,2.55);
\draw[line width=2pt,blue_theme] (5.8,3)--(5.8,3.55) node[above]{$Y$};
\draw[line width=1.6pt,lacuna] (2.65,1.97)--(2.65,4.7)--(8.6,4.7)--(8.6,1.97);
\node[above,lacuna] at (5.6,4.7){entrada $A$};
\node[font=\scriptsize,anchor=west] at (0,2.55){\strut};
\fi
\draw[blue_theme] (0,0) rectangle (11,1.4);
\end{tikzpicture}}

\newcommand{\cicloFabricacao}{%
\begin{tikzpicture}[x=1cm,y=1cm,font=\small]
\foreach \x in {0,3.2,6.4}{\fill[blue_theme!15] (\x,1.5) rectangle ++(2.2,.55);}
\fill[blue_theme!65] (0,2.05) rectangle (2.2,2.4);
\fill[blue_theme!65] (3.2,2.05) rectangle (5.4,2.4);
\foreach \x in {3.2,4.65}{\fill[lacuna!60] (\x,2.4) rectangle ++(.75,.3);}
\foreach \x in {6.4,7.85}{\fill[blue_theme!65] (\x,2.05) rectangle ++(.75,.35);}
\node[font=\bfseries,blue_theme] at (1.1,3.45){Criar uma camada};
\node[font=\bfseries,blue_theme] at (4.3,3.45){Definir um padrão};
\node[font=\bfseries,blue_theme] at (7.5,3.45){Modificar o material};
\node[align=center,font=\scriptsize] at (1.1,.7){Crescer ou depositar\\isolantes e condutores};
\node[align=center,font=\scriptsize] at (4.3,.7){Litografia no resiste\\e alinhamento};
\node[align=center,font=\scriptsize] at (7.5,.7){Gravar, dopar ou abrir\\acesso às camadas};
\draw[seta] (2.4,2.15)--(3,2.15);\draw[seta] (5.6,2.15)--(6.2,2.15);
\draw[seta] (7.5,.1)--(7.5,-.6)--(1.1,-.6)--(1.1,.1);
\node[fill=white,inner sep=2pt,font=\scriptsize] at (4.3,-.6){Repetir com novas máscaras e materiais};
\path[use as bounding box] (-.5,-1) rectangle (9,4);
\end{tikzpicture}}

\newcommand{\escalaAtomica}{%
\begin{tikzpicture}[x=1cm,y=1cm,font=\small]
\node[blue_theme,font=\bfseries] at (1.9,4.1){Tamanho de referência};
\shade[inner color=blue_theme!5,outer color=blue_theme!35] (1.9,1.9) circle(.85);
\node[blue_theme,font=\large] at (1.9,1.9){Si};
\draw[<->,blue_theme] (1.05,.65)--node[below,font=\scriptsize]{diâmetro covalente $\approx0{,}228$ nm}(2.75,.65);
\node[blue_theme,font=\bfseries] at (7.2,4.1){Distância na rede cristalina};
\draw[blue_theme!40,line width=2pt] (6,1.9)--(8.35,1.9);
\foreach \x in {6,8.35}{\shade[inner color=blue_theme!5,outer color=blue_theme!35] (\x,1.9) circle(1.14);\node[blue_theme] at (\x,1.9){Si};}
\draw[densely dashed,blue_theme!50] (6,1.9)--(6,3.45);\draw[densely dashed,blue_theme!50] (8.35,1.9)--(8.35,3.45);
\draw[<->,blue_theme] (6,3.35)--node[above,font=\scriptsize]{$\approx0{,}235$ nm}(8.35,3.35);
\node[align=center,font=\scriptsize] at (7.2,.3){Separação entre vizinhos mais próximos.\\A organização real é tridimensional.};
\path[use as bounding box] (-.2,-.4) rectangle (10,4.6);
\end{tikzpicture}}

\newcommand{\reguaAtomica}{%
\begin{tikzpicture}[x=.74cm,y=1cm,font=\small]
\node[anchor=west,blue_theme,font=\bfseries] at (0,4.3){Comprimentos físicos em uma mesma escala};
\draw[<->,lacuna,thick] (0,3.5)--node[above,font=\scriptsize]{13,5 nm · comprimento de onda EUV}(13.5,3.5);
\draw[<->,blue_theme,thick] (0,2.5)--node[above,font=\scriptsize]{8 nm · resolução nominal High-NA}(8,2.5);
\draw[blue_theme!35,line width=.8pt] (0,1.5)--(8,1.5);
\foreach \i in {0,...,34}{\fill[blue_theme!65] ({\i*.235},1.5) circle[x radius=.095,y radius=.07];}
\node[anchor=west,font=\scriptsize] at (0,.9){8 nm $\approx$ 34 separações de 0,235 nm};
\node[anchor=west,align=left,font=\scriptsize,text width=3.1cm] at (9.1,1.9){\textbf{``2 nm'' é o nome do nó.}\\Não é uma medida física única do transistor.};
\draw[blue_theme!30] (8.65,.7)--(8.65,2.8);
\node[anchor=west,font=\scriptsize] at (0,-.05){Régua unidimensional ilustrativa; não representa a estrutura de um transistor.};
\path[use as bounding box] (-.4,-.4) rectangle (14,4.7);
\end{tikzpicture}}

\newcommand{\barreiraTunelamento}[1]{%
\begin{tikzpicture}[x=1cm,y=1cm,font=\small]
\def\largura{2.7}\ifnum#1=2\def\largura{.85}\fi
\draw[-{Latex},blue_theme] (.2,.3)--(.2,4) node[above,font=\scriptsize]{energia};
\fill[blue_theme!12] (3.4,.55) rectangle (3.4+\largura,3.2);
\draw[blue_theme,very thick] (.5,.55)--(3.4,.55)--(3.4,3.2)--(3.4+\largura,3.2)--(3.4+\largura,.55)--(10.1,.55);
\draw[lacuna,dashed] (.5,1.6)--(10.1,1.6);
\node[anchor=west,font=\scriptsize,lacuna] at (.7,1.95){energia do elétron};
\draw[-{Latex},eletron,very thick] (.8,1.6)--(2.9,1.6);
\fill[eletron] (1.1,1.6) circle(.09);
\draw[<->,blue_theme] (3.4,3.65)--node[above,font=\scriptsize]{espessura da barreira}(3.4+\largura,3.65);
\ifnum#1=1
\draw[-{Latex},eletron!35,dashed,line width=.6pt] (6.3,1.6)--(7.5,1.6);
\node[font=\scriptsize,align=center] at (8,2.65){Transmissão menor\\para esta barreira};
\else
\draw[-{Latex},red,very thick] (4.6,1.6)--(8,1.6);
\node[font=\scriptsize,align=center] at (7.3,2.65){Maior probabilidade\\de atravessar};
\fi
\node[font=\scriptsize,anchor=west] at (.7,-.1){Diagrama de energia: a altura representa energia, não uma parede física.};
\path[use as bounding box] (-.2,-.55) rectangle (10.5,4.35);
\end{tikzpicture}}

% Analogia urbana original: ilustra conectividade; não fixa custos ou ganhos.
\newcommand{\arvoreUrbana}[2]{%
\draw[brown!70!black,line width=1pt] (#1,#2)--++(0,.4);
\fill[green_ufc!45!black] (#1,#2+.5) circle(.17);
\fill[green_ufc!65] (#1-.1,#2+.48) circle(.14);}
\newcommand{\predioUrbano}[5]{%
\begin{scope}[shift={(#1,#2)}]
\fill[black!8] (.12,-.08)--(#3+.52,.1)--(#3+.7,.35)--(.4,.2)--cycle;
\fill[#5!18] (0,0) rectangle (#3,#4);
\fill[#5!38] (#3,0)--(#3+.35,.25)--(#3+.35,#4+.25)--(#3,#4)--cycle;
\fill[#5!8] (0,#4)--(.35,#4+.25)--(#3+.35,#4+.25)--(#3,#4)--cycle;
\draw[#5!65!black,line width=.5pt] (0,0)--(0,#4)--(.35,#4+.25)--(#3+.35,#4+.25)--(#3+.35,.25)--(#3,0)--cycle;
\draw[#5!60!black] (0,#4)--(#3,#4)--(#3,0) (#3,#4)--(#3+.35,#4+.25);
\foreach \yy in {.45,.9,1.35}{\ifdim\yy pt<#4pt\relax\foreach \xx in {.2,.6}{\ifdim\xx pt<#3pt\relax\fill[white!90!#5] (\xx,\yy) rectangle ++(.2,.22);\fi}\fi}
\fill[#5!60!black] (.36,0) rectangle (.61,.34);
\end{scope}}
\newcommand{\pessoaDados}[2]{%
\begin{scope}[shift={(#1,#2)}]
\fill[blue_theme] (0,.27) circle(.065);
\draw[blue_theme,line width=1pt] (0,.2)--(0,.03)--(-.1,-.13) (0,.03)--(.1,-.13) (-.12,.13)--(0,.17)--(.17,.14);
\fill[lacuna!75] (.13,.02) rectangle (.33,.22);\draw[lacuna!70!black] (.13,.02) rectangle (.33,.22);
\end{scope}}
\newcommand{\cidadeDados}[1]{%
\begin{tikzpicture}[x=1cm,y=1cm,font=\small,line join=round]
\ifnum#1=1
% Terreno, calçadas e praças delimitam bairros, sem cartões de texto.
\fill[green_ufc!5] (-.1,.45) rectangle (11.7,4.45);
\fill[blue_theme!6] (.15,1.4) rectangle (3.65,4.2);
\fill[ativo!8] (7.75,1.4) rectangle (11.4,4.2);
\fill[gray!15] (3.8,2.5) rectangle (7.5,4.1);
\foreach \x/\y in {4.1/2.8,5/2.8,5.9/2.8,6.8/2.8,4.1/3.6,6.8/3.6}{\arvoreUrbana{\x}{\y}}
% Via sinuosa: conserva o significado de distância e cruzamentos.
\draw[gray!55,line width=14pt] (2.1,1.8)--(3.5,1.8)--(3.5,.9)--(7.9,.9)--(7.9,1.8)--(9.3,1.8);
\draw[white,line width=.65pt,dash pattern=on 4pt off 3pt] (2.1,1.8)--(3.5,1.8)--(3.5,.9)--(7.9,.9)--(7.9,1.8)--(9.3,1.8);
\foreach \x in {3.32,3.5,3.68}{\draw[white,line width=1.4pt] (\x,1.56)--(\x,2.04);}
\foreach \x in {7.72,7.9,8.08}{\draw[white,line width=1.4pt] (\x,1.56)--(\x,2.04);}
% Edifícios e depósitos do bairro de memória.
\predioUrbano{.55}{2.15}{1.15}{1.8}{blue_theme}
\predioUrbano{2.1}{2.15}{.95}{1.35}{blue_theme}
\predioUrbano{8.45}{2.15}{1.15}{1.8}{ativo}
\predioUrbano{10}{2.15}{.95}{1.35}{ativo}
\foreach \x/\y in {.35/1.8,3.2/3.7,8.25/3.7,11.1/1.8}{\arvoreUrbana{\x}{\y}}
% Caminhão de dados com cabine, carga e rodas.
\begin{scope}[shift={(5.2,.96)}]
\fill[blue_theme] (-.55,-.02) rectangle (.3,.43);\fill[blue_theme!70] (.3,-.02)--(.7,-.02)--(.7,.24)--(.52,.43)--(.3,.43)--cycle;
\fill[white] (.38,.24) rectangle (.57,.34);
\foreach \x in {-.35,.48}{\fill[black!75] (\x,-.08) circle(.1);\fill[white] (\x,-.08) circle(.035);}
\node[font=\tiny,text=white] at (-.12,.2){0101};
\end{scope}
\draw[-{Latex},blue_theme,very thick] (4.1,.43)--(7.2,.43);
\node[font=\scriptsize,blue_theme,fill=white,inner sep=2pt] at (5.7,.05){Dados atravessam a cidade};
\node[blue_theme,font=\bfseries] at (1.9,4.9){Bairro 1 · memória};
\node[ativo,font=\bfseries] at (9.6,4.9){Bairro 2 · processamento};
\node[font=\scriptsize,black!65] at (5.7,4.45){Mais percurso e interfaces de comunicação};
\else
% Corte de um prédio real: fachada lateral, pilares, lajes e três pavimentos.
\fill[black!8] (2.5,-.15)--(9.2,-.15)--(10.3,.3)--(3.6,.3)--cycle;
\fill[blue_theme!18] (8.55,.2)--(9.2,.65)--(9.2,5.1)--(8.55,4.65)--cycle;
\foreach \y in {.5,1.8,3.1}{\foreach \x in {8.67,8.92}{\fill[white!85!blue_theme] (\x,\y+.15)--(\x+.13,\y+.24)--(\x+.13,\y+.73)--(\x,\y+.64)--cycle;}}
\fill[blue_theme!6] (2.8,.25) rectangle (8.5,4.65);
\fill[ativo!8] (2.85,1.65) rectangle (7.5,3.05);
\fill[blue_theme!10] (2.85,3.1) rectangle (7.5,4.5);
\foreach \y in {.2,1.6,3.05,4.5}{%
\fill[gray!45] (2.65,\y) rectangle (8.55,\y+.13);
\fill[gray!22] (2.65,\y+.13)--(3.3,\y+.55)--(9.2,\y+.55)--(8.55,\y+.13)--cycle;}
% Fundo dos andares desenhado à frente apenas com os elementos internos.
\foreach \x in {2.7,7.5,8.4}{\fill[gray!45] (\x,.2) rectangle (\x+.11,4.65);}
% Memória: três estantes com módulos, não janelas abstratas.
\foreach \x in {3.15,4.5,5.85}{%
\fill[blue_theme!60] (\x,3.28) rectangle ++(.9,.95);
\foreach \y in {3.4,3.66,3.92}{\fill[white!85!blue_theme] (\x+.12,\y) rectangle ++(.66,.13);\fill[lacuna] (\x+.65,\y+.035) circle(.025);}}
% Processamento: mesas, telas e profissionais.
\foreach \x in {3.35,5.15}{%
\draw[ativo!70!black,line width=1pt] (\x,2.12)--++(1.15,0) (\x+.1,2.12)--++(0,-.32) (\x+1.05,2.12)--++(0,-.32);
\fill[ativo!60] (\x+.25,2.3) rectangle ++(.6,.37);\fill[white] (\x+.31,2.37) rectangle ++(.48,.23);
\draw[ativo] (\x+.55,2.3)--(\x+.55,2.12);
\fill[blue_theme] (\x+.1,2.65) circle(.075);\draw[blue_theme,thick] (\x+.1,2.56)--(\x+.1,2.27)--(\x+.32,2.21);}
\pessoaDados{7.05}{3.4}
% Térreo: energia e infraestrutura.
\fill[lacuna!70] (3.2,.5) rectangle (4.2,1.25);
\node[font=\large,blue_theme] at (3.7,.88){$+$};
\draw[blue_theme,line width=1pt] (4.45,.6)--(6.9,.6);
\foreach \x in {4.6,5.2,5.8,6.4}{\draw[blue_theme] (\x,.6)--(\x,1.05);}
\node[anchor=east,blue_theme,font=\scriptsize,align=right] at (2.4,3.75){2º andar\\\textbf{Memória}};
\node[anchor=east,ativo,font=\scriptsize,align=right] at (2.4,2.3){1º andar\\\textbf{Processamento}};
\node[anchor=east,black!65,font=\scriptsize,align=right] at (2.4,.85){Térreo\\Alimentação};
% Elevador de dados em uma coluna contínua.
\fill[white!60!blue_theme] (7.65,.35) rectangle (8.25,4.4);
\draw[blue_theme!40] (7.65,.35) rectangle (8.25,4.4);
\draw[<->,blue_theme,very thick] (7.95,1.95)--(7.95,3.9);
\fill[lacuna!75] (7.76,2.9) rectangle (8.14,3.24);\node[font=\tiny,blue_theme] at (7.95,3.07){01};
\node[font=\scriptsize,align=left,anchor=west] at (9.45,2.9){Conexões\\verticais:\\menos percurso,\\mais paralelismo};
\node[font=\scriptsize,blue_theme] at (5.6,5.18){Mesmo prédio: funções em andares especializados};
\fi
\path[use as bounding box] (-.3,-.3) rectangle (12,5.55);
\end{tikzpicture}}

\newcommand{\figura}[1]{\begin{center}\begin{adjustbox}{max width=\linewidth,max totalheight=5.1cm}#1\end{adjustbox}\end{center}}
\author[L. Nogueira / L. Soares]{\texorpdfstring{Luan F. Nogueira\\Luiz F. de Melo Soares}{Luan F. Nogueira e Luiz F. de Melo Soares}}
\institute[UFC]{\department{Minicurso de Microeletrônica}\\\ufc}
\date{SEELECT}

\usepackage{multicol}
% Gerado por scripts/gerar_brasil.py.
% Contorno simplificado: johan/world.geo.json, countries/BRA.geo.json.
% Projeção esquemática linear; mapa de localização por polos, sem escala cartográfica.
\newcommand{\contornoBrasil}{(2.7019,0.6243)--(2.9220,0.8493)--(3.1082,1.0095)--(3.2190,1.0767)--(3.3580,1.1676)--(3.3613,1.2994)--(3.2785,1.3946)--(3.1968,1.3630)--(3.2292,1.4582)--(3.2516,1.5558)--(3.2517,1.6465)--(3.1923,1.6765)--(3.1304,1.6498)--(3.0689,1.6571)--(3.0496,1.7206)--(3.0342,1.8718)--(3.0033,1.9211)--(2.8919,1.9658)--(2.8246,1.9334)--(2.6504,1.9651)--(2.6613,2.1891)--(2.6125,2.2808)--(2.6641,2.3150)--(2.6483,2.4090)--(2.6935,2.4813)--(2.7228,2.6113)--(2.6838,2.7138)--(2.5937,2.7602)--(2.5760,2.8253)--(2.6002,2.9206)--(2.2839,2.9274)--(2.2204,3.1195)--(2.2686,3.1223)--(2.2664,3.1934)--(2.2342,3.2416)--(2.2270,3.3370)--(2.1311,3.3859)--(2.0273,3.3843)--(1.9590,3.4322)--(1.8475,3.4649)--(1.7826,3.5265)--(1.5978,3.5538)--(1.4186,3.7016)--(1.4319,3.8122)--(1.4116,3.8756)--(1.4292,3.9993)--(1.2133,3.9713)--(1.1263,3.9094)--(0.9820,3.8425)--(0.9452,3.7926)--(0.8603,3.7890)--(0.7376,3.8030)--(0.6445,3.7745)--(0.5695,3.7935)--(0.5805,4.0441)--(0.4451,3.9469)--(0.2995,3.9512)--(0.2371,4.0392)--(0.1276,4.0487)--(0.1625,4.1196)--(0.0708,4.2200)--(0.0021,4.3686)--(0.0456,4.3987)--(0.0455,4.4684)--(0.1452,4.5161)--(0.1287,4.6053)--(0.1709,4.6627)--(0.1828,4.7397)--(0.3715,4.8520)--(0.5067,4.8837)--(0.5289,4.9085)--(0.6776,4.9008)--(0.7517,5.3532)--(0.7556,5.4248)--(0.7298,5.5193)--(0.6566,5.5794)--(0.6574,5.6993)--(0.7504,5.7265)--(0.7833,5.7094)--(0.7889,5.7726)--(0.6922,5.7897)--(0.6902,5.8929)--(1.0117,5.8893)--(1.0663,5.9461)--(1.1121,5.8938)--(1.1443,5.7965)--(1.1754,5.8168)--(1.2662,5.7295)--(1.3945,5.7402)--(1.4265,5.7907)--(1.5492,5.8292)--(1.6171,5.8563)--(1.6363,5.9262)--(1.7541,5.9731)--(1.7452,6.0078)--(1.6055,6.0220)--(1.5825,6.1259)--(1.5892,6.2365)--(1.5153,6.2793)--(1.5463,6.2945)--(1.6684,6.2734)--(1.7996,6.2321)--(1.8473,6.2711)--(1.9659,6.2968)--(2.1505,6.3585)--(2.2108,6.4215)--(2.1890,6.4680)--(2.2747,6.4753)--(2.3131,6.4373)--(2.2917,6.3649)--(2.3484,6.3399)--(2.3862,6.2632)--(2.3405,6.2051)--(2.3142,6.0646)--(2.3564,5.9812)--(2.3684,5.9048)--(2.4699,5.8274)--(2.5509,5.8192)--(2.5691,5.8516)--(2.6213,5.8587)--(2.6959,5.8876)--(2.7496,5.9315)--(2.8409,5.9175)--(2.8810,5.9234)--(2.9707,5.9099)--(2.9856,5.9436)--(2.9579,5.9764)--(2.9744,6.0242)--(3.0410,6.0095)--(3.1189,6.0264)--(3.2134,5.9915)--(3.2855,5.9574)--(3.3365,6.0022)--(3.3735,5.9953)--(3.3960,5.9488)--(3.4750,5.9606)--(3.5382,6.0233)--(3.5889,6.1448)--(3.6865,6.2958)--(3.7427,6.3036)--(3.7835,6.2123)--(3.8760,5.9238)--(3.9643,5.8965)--(3.9687,5.7826)--(3.8446,5.6468)--(3.8959,5.5971)--(4.1876,5.5711)--(4.1936,5.4058)--(4.3189,5.5140)--(4.5265,5.4547)--(4.8006,5.3540)--(4.8811,5.2573)--(4.8540,5.1659)--(5.0459,5.2168)--(5.3670,5.1295)--(5.6135,5.1359)--(5.8574,4.9994)--(6.0682,4.8145)--(6.1953,4.7669)--(6.3364,4.7603)--(6.3962,4.7083)--(6.4522,4.4982)--(6.4796,4.3984)--(6.4138,4.1256)--(6.3299,4.0179)--(6.0973,3.7883)--(5.9922,3.6018)--(5.8701,3.4587)--(5.8288,3.4555)--(5.7827,3.3341)--(5.7944,3.0249)--(5.7484,2.7706)--(5.7309,2.6618)--(5.6787,2.5967)--(5.6495,2.3761)--(5.4822,2.1608)--(5.4541,1.9903)--(5.3206,1.9188)--(5.2819,1.8199)--(5.1027,1.8203)--(4.8431,1.7569)--(4.7269,1.6835)--(4.5421,1.6353)--(4.3479,1.5039)--(4.2082,1.3403)--(4.1842,1.2171)--(4.2117,1.1260)--(4.1808,0.9593)--(4.1434,0.8788)--(4.0281,0.7880)--(3.8450,0.4976)--(3.6999,0.3667)--(3.5877,0.2895)--(3.5125,0.1326)--(3.4033,0.0382)--(3.3577,0.1317)--(3.4304,0.2099)--(3.3350,0.3222)--(3.2055,0.4134)--(3.0358,0.5191)--(2.9744,0.5143)--(2.8090,0.6419)--(2.7019,0.6243)--cycle}

\newcommand{\mapaBrasil}[1]{%
\begin{tikzpicture}[x=.75cm,y=.75cm,font=\scriptsize]
\ifdefined\SemRotulos\path[use as bounding box](-.3,-.6) rectangle (6.8,7);\else\path[use as bounding box](-.3,-.6) rectangle (11.7,7);\fi
\filldraw[fill=blue_theme!5,draw=blue_theme,line width=.8pt] \contornoBrasil;
\fill[black!22] (2.307,5.098)circle(.05);
\fill[black!22] (4.310,3.005)circle(.05);
\fill[black!22] (6.290,4.417)circle(.05);
\fill[black!22] (5.857,3.470)circle(.05);
\fill[black!22] (4.637,2.062)circle(.05);
\fill[black!22] (4.960,2.323)circle(.05);
\fill[black!22] (4.445,1.832)circle(.05);
\fill[black!22] (4.199,1.058)circle(.05);
\fill[black!22] (3.757,0.655)circle(.05);
\fill[black!22] (4.059,2.062)circle(.05);
\fill[black!22] (4.080,1.414)circle(.05);
\fill[black!22] (3.218,0.693)circle(.05);
\fill[black!22] (4.640,1.785)circle(.05);
\foreach \selecionado in {#1}{
\ifnum\selecionado=1\relax\ifdefined\SemRotulos\else\draw[green_ufc,line width=.6pt](2.307,5.098)--(1.50,5.90);\node[anchor=west,fill=white,inner sep=1.5pt,text=green_ufc,font=\scriptsize] at(1.50,5.90){Manaus};\fi\node[circle,fill=green_ufc,text=white,inner sep=1.3pt,font=\tiny] at(2.307,5.098){1};\fi
\ifnum\selecionado=2\relax\ifdefined\SemRotulos\else\draw[red!65!blue,line width=.6pt](4.310,3.005)--(2.80,3.80);\node[anchor=west,fill=white,inner sep=1.5pt,text=red!65!blue,font=\scriptsize] at(2.80,3.80){Brasília};\fi\node[circle,fill=red!65!blue,text=white,inner sep=1.3pt,font=\tiny] at(4.310,3.005){2};\fi
\ifnum\selecionado=3\relax\ifdefined\SemRotulos\else\draw[lacuna,line width=.6pt](6.290,4.417)--(7.50,5.10);\node[anchor=west,fill=white,inner sep=1.5pt,text=lacuna,font=\scriptsize] at(7.50,5.10){Campina Grande};\fi\node[circle,fill=lacuna,text=white,inner sep=1.3pt,font=\tiny] at(6.290,4.417){3};\fi
\ifnum\selecionado=4\relax\ifdefined\SemRotulos\else\draw[lacuna,line width=.6pt](5.857,3.470)--(7.50,3.80);\node[anchor=west,fill=white,inner sep=1.5pt,text=lacuna,font=\scriptsize] at(7.50,3.80){Salvador};\fi\node[circle,fill=lacuna,text=white,inner sep=1.3pt,font=\tiny] at(5.857,3.470){4};\fi
\ifnum\selecionado=5\relax\ifdefined\SemRotulos\else\draw[blue_theme,line width=.6pt](4.837,2.143)--(7.50,2.90);\node[anchor=west,fill=white,inner sep=1.5pt,text=blue_theme,font=\scriptsize] at(7.50,2.90){Interior de MG};\fi\node[circle,fill=blue_theme,text=white,inner sep=1.3pt,font=\tiny] at(4.837,2.143){5};\fi
\ifnum\selecionado=6\relax\ifdefined\SemRotulos\else\draw[blue_theme,line width=.6pt](5.010,2.473)--(7.50,3.35);\node[anchor=west,fill=white,inner sep=1.5pt,text=blue_theme,font=\scriptsize] at(7.50,3.35){Belo Horizonte};\fi\node[circle,fill=blue_theme,text=white,inner sep=1.3pt,font=\tiny] at(5.010,2.473){6};\fi
\ifnum\selecionado=7\relax\ifdefined\SemRotulos\else\draw[blue_theme,line width=.6pt](4.365,1.932)--(7.50,2.10);\node[anchor=west,fill=white,inner sep=1.5pt,text=blue_theme,font=\scriptsize] at(7.50,2.10){Campinas};\fi\node[circle,fill=blue_theme,text=white,inner sep=1.3pt,font=\tiny] at(4.365,1.932){7};\fi
\ifnum\selecionado=8\relax\ifdefined\SemRotulos\else\draw[ativo,line width=.6pt](4.299,1.078)--(7.50,0.90);\node[anchor=west,fill=white,inner sep=1.5pt,text=ativo,font=\scriptsize] at(7.50,0.90){Florianópolis};\fi\node[circle,fill=ativo,text=white,inner sep=1.3pt,font=\tiny] at(4.299,1.078){8};\fi
\ifnum\selecionado=9\relax\ifdefined\SemRotulos\else\draw[ativo,line width=.6pt](3.757,0.655)--(5.00,-0.25);\node[anchor=west,fill=white,inner sep=1.5pt,text=ativo,font=\scriptsize] at(5.00,-0.25){Porto Alegre};\fi\node[circle,fill=ativo,text=white,inner sep=1.3pt,font=\tiny] at(3.757,0.655){9};\fi
\ifnum\selecionado=10\relax\ifdefined\SemRotulos\else\draw[blue_theme,line width=.6pt](3.919,2.163)--(1.40,2.50);\node[anchor=west,fill=white,inner sep=1.5pt,text=blue_theme,font=\scriptsize] at(1.40,2.50){Interior de SP};\fi\node[circle,fill=blue_theme,text=white,inner sep=1.3pt,font=\tiny] at(3.919,2.163){10};\fi
\ifnum\selecionado=11\relax\ifdefined\SemRotulos\else\draw[ativo,line width=.6pt](4.050,1.384)--(7.50,1.45);\node[anchor=west,fill=white,inner sep=1.5pt,text=ativo,font=\scriptsize] at(7.50,1.45){Curitiba};\fi\node[circle,fill=ativo,text=white,inner sep=1.3pt,font=\tiny] at(4.050,1.384){11};\fi
\ifnum\selecionado=12\relax\ifdefined\SemRotulos\else\draw[ativo,line width=.6pt](3.218,0.693)--(1.50,0.25);\node[anchor=west,fill=white,inner sep=1.5pt,text=ativo,font=\scriptsize] at(1.50,0.25){Interior do RS};\fi\node[circle,fill=ativo,text=white,inner sep=1.3pt,font=\tiny] at(3.218,0.693){12};\fi
\ifnum\selecionado=13\relax\ifdefined\SemRotulos\else\draw[blue_theme,line width=.6pt](4.840,1.685)--(7.50,2.55);\node[anchor=west,fill=white,inner sep=1.5pt,text=blue_theme,font=\scriptsize] at(7.50,2.55){S. José dos Campos};\fi\node[circle,fill=blue_theme,text=white,inner sep=1.3pt,font=\tiny] at(4.840,1.685){13};\fi
}
\end{tikzpicture}}
\newcommand{\creditoBrasil}{\fonte{Levantamento: Matheus Grossi · \href{https://github.com/MattGrossi12/brazil_silicon_roadmap}{Brazil Silicon Roadmap}, planilha PT-BR fornecida. Mapa próprio; polos esquemáticos.}}

\title[Aula 3 • Caminhos e Brasil]{\huge\textbf{Do transistor ao chip}}
\begin{document}
\capa{Encerramento da aula 3 · Para onde seguir?}{Trilhas de estudo e microeletrônica no Brasil • cerca de 25 minutos}
\note{Bloco final da oficina de duas horas. Reservar cerca de 95 minutos à futura prática e 25 ao encerramento; a prática continua sem roteiro elaborado. As páginas regionais do Brasil são consultáveis: destacar exemplos, sem ler toda a relação.}

\section{Trilhas}
\begin{frame}{Uma base, cinco caminhos que se conectam}
\figura{\begin{tikzpicture}[x=1cm,y=1cm,font=\small]
\node[circle,draw=blue_theme,fill=white,minimum size=22mm,align=center](base)at(0,0){MOSFET\\e CMOS};
\foreach \y/\t/\ex in {2.4/Digital/Portas e RTL,1.2/Arquitetura/Processadores,0/Analógica/Amplificadores,-1.2/Sinais mistos/Conversores,-2.4/RF/Comunicação}{\node[anchor=west,align=left] (p\y) at(3.2,\y){\textbf{\t}\\\ex};\draw[seta](base.east)--(3,\y);}
\end{tikzpicture}}
\chamada{Microeletrônica reúne essas áreas e outras especialidades.}
\end{frame}
\note{O objetivo é abrir possibilidades, não exigir escolha de carreira. Arquitetura conecta escolhas do sistema a implementações digitais; RF se apoia em analógica, sinais e eletromagnetismo. Nenhuma trilha é totalmente isolada.}

\begin{frame}{Digital: implementar funções em hardware}
\figura{\begin{tikzpicture}[node distance=8mm,font=\small]
\node[draw=blue_theme,minimum width=2cm,minimum height=1cm](a){Portas};\node[draw=blue_theme,minimum width=2cm,minimum height=1cm,right=of a](b){Registradores};\node[draw=blue_theme,minimum width=2cm,minimum height=1cm,right=of b](c){ULA};\draw[seta](a)--(b);\draw[seta](b)--(c);\end{tikzpicture}}
\begin{itemize}\item Estudar: lógica sequencial, HDL/RTL e verificação.
\item Primeiro projeto: ULA de 4 bits com testes automáticos.\end{itemize}
\end{frame}
\note{HDL: linguagem de descrição de hardware. RTL: descrição no nível de transferências entre registradores. ULA: unidade lógica e aritmética. As setas mostram um percurso de estudo, não a conectividade interna de uma ULA. FPGA é uma possibilidade de prototipagem; ASIC tem outras etapas de implementação física.}

\begin{frame}{Arquitetura: organizar a computação}
\figura{\begin{tikzpicture}[node distance=12mm,font=\small]
\node[draw=blue_theme,minimum width=2cm,minimum height=2cm,align=center](r){Banco de\\registradores};
\node[draw=blue_theme,minimum width=1.7cm,minimum height=1.2cm,right=of r](u){ULA};\node[draw=blue_theme,minimum width=2cm,minimum height=2cm,right=of u](m){Memória};
\draw[seta](r)--(u);\draw[seta](u)--(m);\draw[seta](m.south)--++(0,-.7)-|(r.south);
\node[above=8mm of u](c){Controle e instruções};\draw[seta](c)--(u);
\end{tikzpicture}}
\begin{itemize}\item Estudar: instruções, caminho de dados e memória.
\item Primeiro projeto: executar uma tarefa e contar ciclos.\end{itemize}
\end{frame}
\note{Diagrama funcional simplificado. A ISA é o contrato de instruções visível ao software; a microarquitetura organiza sua execução. Depois vêm temas como pipeline, caches e paralelismo. O trabalho de Luan com aceleradores ilustra a ligação com essa trilha.}

\begin{frame}{Analógica: trabalhar com sinais contínuos}
\figura{\begin{tikzpicture}[x=1cm,y=1cm,font=\small]
\draw[blue_theme,thick,domain=0:2,samples=70]plot(\x,{.3*sin(360*\x)});
\draw[seta](2.4,0)--(3,0);\draw[blue_theme,thick](3.3,-1)--(3.3,1)--(5.2,0)--cycle;\node at(3.6,.45){$+$};\node at(3.6,-.45){$-$};\draw[seta](5.4,0)--(6,0);
\draw[ativo,thick,domain=6.3:8.3,samples=70]plot(\x,{.9*sin(360*\x)});
\end{tikzpicture}}
\begin{itemize}\item Estudar: polarização, pequenos sinais e realimentação.
\item Primeiro projeto: amplificador fonte comum.\end{itemize}
\end{frame}
\note{O desenho ilustra amplificação ideal de uma senoide, sem representar uma topologia ou ganho medido. O mesmo MOSFET passa a ser visto como dispositivo de controle contínuo de corrente. Projetos reais também precisam de ruído, estabilidade, faixa de saída e variações. O trabalho de Luiz com amplificadores e layout é um exemplo.}

\begin{frame}{Sinais mistos: do sinal ao código}
\figura{\begin{tikzpicture}[x=1cm,y=1cm,font=\small]
\draw[blue_theme,thick,domain=0:2.5,samples=70]plot(\x,{.7*sin(180*\x)});
\draw[seta](2.9,0)--(3.5,0);\node[draw=blue_theme,minimum width=1.7cm,minimum height=1.3cm](c)at(4.6,0){ADC};\draw[seta](5.6,0)--(6.2,0);
\node[align=left,font=\ttfamily\small] at(7.1,0){001\\011\\110\\100};
\node at(1.2,-1.3){Tensão};\node at(4.6,-1.3){Amostragem};\node at(7.1,-1.3){Bits};
\end{tikzpicture}}
\begin{itemize}\item Estudar: amostragem, quantização e conversores.
\item Primeiro projeto: modelo de ADC de 3 bits.\end{itemize}
\end{frame}
\note{ADC: conversor analógico--digital. DAC faz a conversão no sentido oposto. O código mostrado é ilustrativo, sem correspondência numérica com as amostras desenhadas. A área também envolve interfaces de sensores, sincronização e convivência entre blocos analógicos e digitais.}

\begin{frame}{RF integrada: comunicar por ondas}
\figura{\begin{tikzpicture}[x=1cm,y=1cm,font=\small]
\draw[blue_theme,line width=1.5pt](0,-1)--(0,1.2);\draw[blue_theme,line width=1.5pt](-.6,1.2)--(0,.6)--(.6,1.2);
\draw[blue_theme](-.8,.2)arc(210:150:1);\draw[blue_theme](.8,1.2)arc(30:-30:1);
\draw[seta](.5,-.4)--(1.5,-.4);\node[draw=blue_theme,minimum width=2cm,minimum height=1.1cm]at(2.7,-.4){Amplificador};\draw[seta](3.9,-.4)--(4.7,-.4);\node[draw=blue_theme,minimum width=2cm,minimum height=1.1cm]at(5.9,-.4){Misturador};\draw[seta](7.1,-.4)--(8,-.4);
\end{tikzpicture}}
\begin{itemize}\item Estudar: analógica, sinais, impedâncias e eletromagnetismo.
\item Primeiro projeto: rede LC de casamento de impedâncias.\end{itemize}
\end{frame}
\note{LC: indutor e capacitor. A cadeia é uma ilustração funcional de recepção, não um rádio completo. Ruído, linearidade, componentes parasitas e encapsulamento são importantes. Não definir a área por uma frequência de corte universal.}

\begin{frame}{As trilhas se encontram em um produto}
\figura{\begin{tikzpicture}[node distance=4mm,font=\small]
\node[draw=blue_theme,align=center,minimum width=1.5cm,minimum height=1cm](s){Sensor};\node[draw=blue_theme,align=center,minimum width=1.8cm,minimum height=1cm,right=of s](a){Analógica};\node[draw=blue_theme,align=center,minimum width=1.5cm,minimum height=1cm,right=of a](d){ADC};\node[draw=blue_theme,align=center,minimum width=1.6cm,minimum height=1cm,right=of d](p){Digital};\node[draw=blue_theme,align=center,minimum width=1.3cm,minimum height=1cm,right=of p](r){RF};\draw[seta](s)--(a);\draw[seta](a)--(d);\draw[seta](d)--(p);\draw[seta](p)--(r);
\end{tikzpicture}}
\chamada{Exemplo: um nó de sensores sem fio.}
\small\centering Projeto físico • verificação • processos • teste • encapsulamento • ferramentas
\end{frame}
\note{O processamento e a comunicação reais incluem interfaces e conversões adicionais. A arquitetura organiza dados, memória e consumo. O setor também precisa de funções transversais: não se resume a desenhar transistores ou escrever HDL. Cada bloco pode estar no mesmo chip ou em componentes separados.}

% Gerado por scripts/gerar_brasil.py. Relação integral de 73 presenças.
\section{Âmbito Brasil}
\begin{frame}{Microeletrônica no Brasil}
\duas{\begingroup\def\SemRotulos{1}\figura{\mapaBrasil{1,2,3,4,8,9}}\endgroup}{\vspace{5mm}\textcolor{blue_theme}{\Huge 54}\ organizações\par\medskip\textcolor{blue_theme}{\Huge 13}\ polos\par\medskip Projeto • ferramentas • integração • fabricação • pesquisa}
\creditoBrasil
\end{frame}
\note{Os números são da planilha local fornecida, com 54 organizações e 73 presenças em 13 polos. Diferem do resumo de outra versão da página do repositório. Um nome pode aparecer em vários polos. O mapa é próprio, com contorno geográfico simplificado e localizações esquemáticas agregadas. As próximas páginas preservam toda a relação; destaque exemplos em vez de ler cada nome. Centros e grupos de pesquisa são distinguidos de empresas. Créditos do levantamento: Matheus Grossi, Brazil Silicon Roadmap.}
\begin{frame}{Âmbito Brasil · Norte e Nordeste}
\begin{columns}[T,onlytextwidth]
\begin{column}{.35\textwidth}
\figura{\mapaBrasil{1,3,4}}
{\tiny $\bullet$ empresa\par $\circ$ centro/grupo}
\end{column}
\begin{column}{.63\textwidth}
\footnotesize\setlength{\columnsep}{3mm}\begin{multicols}{2}\raggedright
\textcolor{blue_theme}{\textbf{Manaus}}\par\smallskip
\noindent\parbox{\linewidth}{$\circ$ Instituto SENAI de Inovação em Microeletrônica}\par
\noindent\parbox{\linewidth}{$\bullet$ ADATA}\par
\noindent\parbox{\linewidth}{$\bullet$ Zilia}\par
\noindent\parbox{\linewidth}{$\bullet$ Stoneridge}\par
\noindent\parbox{\linewidth}{$\circ$ HBR}\par
\noindent\parbox{\linewidth}{$\circ$ Instituto de Pesquisas ELDORADO}\par
\noindent\parbox{\linewidth}{$\bullet$ HT Micron}\par
\noindent\parbox{\linewidth}{$\bullet$ CCBS / TERA}\par
\medskip
\textcolor{blue_theme}{\textbf{Campina Grande}}\par\smallskip
\noindent\parbox{\linewidth}{$\bullet$ HWiT}\par
\medskip
\textcolor{blue_theme}{\textbf{Salvador}}\par\smallskip
\noindent\parbox{\linewidth}{$\circ$ SENAI CIMATEC}\par
\noindent\parbox{\linewidth}{$\circ$ HBR}\par
\medskip
\end{multicols}
\end{column}
\end{columns}
\creditoBrasil
\end{frame}
\note{Relação da planilha PT-BR fornecida. Os polos abrangem cidades e proximidades ou áreas de interior; os pontos não indicam endereços individuais. As categorias e incertezas de atuação estão preservadas nos CSVs do projeto. Algumas organizações atuam em módulos, ferramentas ou eletrônica aplicada, sem confirmação de projeto local de CI/RTL. Os nomes estão aqui para consulta; escolha um ou dois exemplos para explicar a diversidade da cadeia.}
\begin{frame}{Âmbito Brasil · Centro-Oeste e Minas Gerais}
\begin{columns}[T,onlytextwidth]
\begin{column}{.35\textwidth}
\figura{\mapaBrasil{2,5,6}}
{\tiny $\bullet$ empresa\par $\circ$ centro/grupo}
\end{column}
\begin{column}{.63\textwidth}
\footnotesize\setlength{\columnsep}{3mm}\begin{multicols}{2}\raggedright
\textcolor{blue_theme}{\textbf{Brasília}}\par\smallskip
\noindent\parbox{\linewidth}{$\circ$ HBR}\par
\noindent\parbox{\linewidth}{$\circ$ Instituto de Pesquisas ELDORADO}\par
\medskip
\textcolor{blue_theme}{\textbf{Interior de MG}}\par\smallskip
\noindent\parbox{\linewidth}{$\circ$ LC-e}\par
\noindent\parbox{\linewidth}{$\circ$ Inatel}\par
\noindent\parbox{\linewidth}{$\bullet$ Modirum | GESPI}\par
\noindent\parbox{\linewidth}{$\bullet$ Marelli}\par
\noindent\parbox{\linewidth}{$\bullet$ Stellantis}\par
\noindent\parbox{\linewidth}{$\bullet$ Multilaser}\par
\medskip
\textcolor{blue_theme}{\textbf{Belo Horizonte}}\par\smallskip
\noindent\parbox{\linewidth}{$\bullet$ Cadence}\par
\noindent\parbox{\linewidth}{$\bullet$ RFIDO Design}\par
\medskip
\end{multicols}
\end{column}
\end{columns}
\creditoBrasil
\end{frame}
\note{Relação da planilha PT-BR fornecida. Os polos abrangem cidades e proximidades ou áreas de interior; os pontos não indicam endereços individuais. As categorias e incertezas de atuação estão preservadas nos CSVs do projeto. Algumas organizações atuam em módulos, ferramentas ou eletrônica aplicada, sem confirmação de projeto local de CI/RTL. Os nomes estão aqui para consulta; escolha um ou dois exemplos para explicar a diversidade da cadeia.}
\begin{frame}{Âmbito Brasil · São Paulo · polo de Campinas}
\begin{columns}[T,onlytextwidth]
\begin{column}{.35\textwidth}
\figura{\mapaBrasil{7}}
{\tiny $\bullet$ empresa\par $\circ$ centro/grupo}
\end{column}
\begin{column}{.63\textwidth}
\footnotesize\setlength{\columnsep}{3mm}\begin{multicols}{2}\raggedright
\textcolor{blue_theme}{\textbf{Campinas}}\par\smallskip
\noindent\parbox{\linewidth}{$\circ$ Von Braun Labs}\par
\noindent\parbox{\linewidth}{$\bullet$ ADATA}\par
\noindent\parbox{\linewidth}{$\circ$ HBR}\par
\noindent\parbox{\linewidth}{$\bullet$ HWiT}\par
\noindent\parbox{\linewidth}{$\bullet$ EnSilica}\par
\noindent\parbox{\linewidth}{$\circ$ CPQD}\par
\noindent\parbox{\linewidth}{$\bullet$ Aedvices}\par
\noindent\parbox{\linewidth}{$\circ$ Instituto de Pesquisas ELDORADO}\par
\noindent\parbox{\linewidth}{$\bullet$ Stoneridge}\par
\noindent\parbox{\linewidth}{$\bullet$ Valeo}\par
\noindent\parbox{\linewidth}{$\bullet$ Cadence}\par
\noindent\parbox{\linewidth}{$\circ$ CNPEM}\par
\noindent\parbox{\linewidth}{$\circ$ CTI Renato Archer}\par
\noindent\parbox{\linewidth}{$\bullet$ HCLTech}\par
\medskip
\end{multicols}
\end{column}
\end{columns}
\creditoBrasil
\end{frame}
\note{Relação da planilha PT-BR fornecida. Os polos abrangem cidades e proximidades ou áreas de interior; os pontos não indicam endereços individuais. As categorias e incertezas de atuação estão preservadas nos CSVs do projeto. Algumas organizações atuam em módulos, ferramentas ou eletrônica aplicada, sem confirmação de projeto local de CI/RTL. Os nomes estão aqui para consulta; escolha um ou dois exemplos para explicar a diversidade da cadeia.}
\begin{frame}{Âmbito Brasil · São Paulo · outros polos}
\begin{columns}[T,onlytextwidth]
\begin{column}{.35\textwidth}
\figura{\mapaBrasil{10,13}}
{\tiny $\bullet$ empresa\par $\circ$ centro/grupo}
\end{column}
\begin{column}{.63\textwidth}
\footnotesize\setlength{\columnsep}{3mm}\begin{multicols}{2}\raggedright
\textcolor{blue_theme}{\textbf{Interior de SP}}\par\smallskip
\noindent\parbox{\linewidth}{$\bullet$ Macnica DHW}\par
\noindent\parbox{\linewidth}{$\bullet$ WG3}\par
\noindent\parbox{\linewidth}{$\circ$ Chip Inventor}\par
\noindent\parbox{\linewidth}{$\bullet$ Beontag}\par
\noindent\parbox{\linewidth}{$\bullet$ SPICTEC}\par
\noindent\parbox{\linewidth}{$\circ$ LSI-TEC}\par
\noindent\parbox{\linewidth}{$\bullet$ Zilia}\par
\medskip
\textcolor{blue_theme}{\textbf{S. José dos Campos}}\par\smallskip
\noindent\parbox{\linewidth}{$\bullet$ Boeing}\par
\noindent\parbox{\linewidth}{$\bullet$ Embraer}\par
\noindent\parbox{\linewidth}{$\bullet$ Thales Group}\par
\noindent\parbox{\linewidth}{$\bullet$ SIATT}\par
\noindent\parbox{\linewidth}{$\bullet$ IACIT}\par
\noindent\parbox{\linewidth}{$\circ$ INPE}\par
\noindent\parbox{\linewidth}{$\bullet$ Modirum | GESPI}\par
\noindent\parbox{\linewidth}{$\bullet$ Tecsys do Brasil}\par
\medskip
\end{multicols}
\end{column}
\end{columns}
\creditoBrasil
\end{frame}
\note{Relação da planilha PT-BR fornecida. Os polos abrangem cidades e proximidades ou áreas de interior; os pontos não indicam endereços individuais. As categorias e incertezas de atuação estão preservadas nos CSVs do projeto. Algumas organizações atuam em módulos, ferramentas ou eletrônica aplicada, sem confirmação de projeto local de CI/RTL. Os nomes estão aqui para consulta; escolha um ou dois exemplos para explicar a diversidade da cadeia.}
\begin{frame}{Âmbito Brasil · Paraná e Santa Catarina}
\begin{columns}[T,onlytextwidth]
\begin{column}{.35\textwidth}
\figura{\mapaBrasil{11,8}}
{\tiny $\bullet$ empresa\par $\circ$ centro/grupo}
\end{column}
\begin{column}{.63\textwidth}
\footnotesize\setlength{\columnsep}{3mm}\begin{multicols}{2}\raggedright
\textcolor{blue_theme}{\textbf{Curitiba}}\par\smallskip
\noindent\parbox{\linewidth}{$\bullet$ HCLTech}\par
\noindent\parbox{\linewidth}{$\bullet$ Thales Group}\par
\noindent\parbox{\linewidth}{$\circ$ Lactec}\par
\medskip
\textcolor{blue_theme}{\textbf{Florianópolis}}\par\smallskip
\noindent\parbox{\linewidth}{$\bullet$ Macnica DHW}\par
\noindent\parbox{\linewidth}{$\bullet$ Chipus}\par
\noindent\parbox{\linewidth}{$\bullet$ Synopsys}\par
\noindent\parbox{\linewidth}{$\circ$ OSIRIS}\par
\noindent\parbox{\linewidth}{$\circ$ Silicium}\par
\medskip
\end{multicols}
\end{column}
\end{columns}
\creditoBrasil
\end{frame}
\note{Relação da planilha PT-BR fornecida. Os polos abrangem cidades e proximidades ou áreas de interior; os pontos não indicam endereços individuais. As categorias e incertezas de atuação estão preservadas nos CSVs do projeto. Algumas organizações atuam em módulos, ferramentas ou eletrônica aplicada, sem confirmação de projeto local de CI/RTL. Os nomes estão aqui para consulta; escolha um ou dois exemplos para explicar a diversidade da cadeia.}
\begin{frame}{Âmbito Brasil · Rio Grande do Sul}
\begin{columns}[T,onlytextwidth]
\begin{column}{.35\textwidth}
\figura{\mapaBrasil{9,12}}
{\tiny $\bullet$ empresa\par $\circ$ centro/grupo}
\end{column}
\begin{column}{.63\textwidth}
\footnotesize\setlength{\columnsep}{3mm}\begin{multicols}{2}\raggedright
\textcolor{blue_theme}{\textbf{Porto Alegre}}\par\smallskip
\noindent\parbox{\linewidth}{$\bullet$ Impinj}\par
\noindent\parbox{\linewidth}{$\bullet$ SILVACO}\par
\noindent\parbox{\linewidth}{$\bullet$ HCLTech}\par
\noindent\parbox{\linewidth}{$\bullet$ AEL Sistemas}\par
\noindent\parbox{\linewidth}{$\bullet$ Grupo Tellescom}\par
\noindent\parbox{\linewidth}{$\bullet$ Valeo}\par
\noindent\parbox{\linewidth}{$\circ$ Instituto de Pesquisas ELDORADO}\par
\noindent\parbox{\linewidth}{$\bullet$ EnSilica}\par
\noindent\parbox{\linewidth}{$\circ$ itt Chip — Unisinos}\par
\noindent\parbox{\linewidth}{$\bullet$ CEITEC}\par
\noindent\parbox{\linewidth}{$\circ$ Instituto SENAI (RS)}\par
\noindent\parbox{\linewidth}{$\bullet$ InPlace}\par
\noindent\parbox{\linewidth}{$\bullet$ HT Micron}\par
\medskip
\textcolor{blue_theme}{\textbf{Interior do RS}}\par\smallskip
\noindent\parbox{\linewidth}{$\bullet$ TUNA Semicondutores}\par
\noindent\parbox{\linewidth}{$\circ$ GMicro UFSM}\par
\medskip
\end{multicols}
\end{column}
\end{columns}
\creditoBrasil
\end{frame}
\note{Relação da planilha PT-BR fornecida. Os polos abrangem cidades e proximidades ou áreas de interior; os pontos não indicam endereços individuais. As categorias e incertezas de atuação estão preservadas nos CSVs do projeto. Algumas organizações atuam em módulos, ferramentas ou eletrônica aplicada, sem confirmação de projeto local de CI/RTL. Os nomes estão aqui para consulta; escolha um ou dois exemplos para explicar a diversidade da cadeia.}


\section{Próximo passo}
\begin{frame}{Escolha uma pergunta e faça um pequeno projeto}
\figura{\begin{tikzpicture}[node distance=10mm,font=\small]
\node[circle,draw=blue_theme,minimum size=15mm](a){Escolher};\node[circle,draw=blue_theme,minimum size=15mm,right=of a](b){Estudar};\node[circle,draw=blue_theme,minimum size=15mm,right=of b](c){Construir};\node[circle,draw=blue_theme,minimum size=15mm,right=of c](d){Explicar};\draw[seta](a)--(b);\draw[seta](b)--(c);\draw[seta](c)--(d);
\end{tikzpicture}}
\chamada{Qual trilha? Qual primeiro projeto? Como demonstrar que funciona?}
\end{frame}
\note{Fechar com dois ou três relatos curtos. A turma pode registrar uma trilha, um pré-requisito e uma entrega. Não é preciso decidir uma carreira hoje. Os mapas e as listas ficam para consulta.}

\begin{frame}{Créditos e caminhos de consulta}
\small\begin{itemize}\item Desenvolvimento: Luan F. Nogueira e Luiz F. de Melo Soares.
\item Levantamento Brasil: \textbf{Matheus Grossi}, \href{https://github.com/MattGrossi12/brazil_silicon_roadmap}{Brazil Silicon Roadmap}; planilha PT-BR fornecida.
\item Contorno: \href{https://github.com/johan/world.geo.json/blob/master/countries/BRA.geo.json}{johan/world.geo.json}; mapa redesenhado em TikZ.
\item Continuação de estudos: \href{https://ocw.mit.edu/courses/6-004-computation-structures-spring-2017/}{MIT 6.004}, Salem e Dutertre.
\item Template Beamer UFC: Maurício Moreira Neto.\end{itemize}
\end{frame}
\note{As localizações são polos agregados e não endereços. As incertezas e categorias da planilha estão preservadas nos dados derivados. O mapa não reproduz a imagem fornecida. As sequências de estudo e os pequenos projetos são propostas didáticas, não promessas de emprego ou ofertas de curso.}
\end{document}
